\section{Aquisição e preparação dos sinais}
\label{sec:aquisicao}

\subsection{Características dos componentes}

A Tabela~\ref{tab:componentes} reúne informações de referência e decisões da
montagem. Especificações de catálogo não equivalem a erros determinados no
protótipo: posicionamento, fiação, alimentação e características do módulo
podem afetar as leituras. Também é necessário distinguir resolução digital,
erro especificado e incerteza de uma comparação experimental.

\begin{table}[htbp]
\centering
\caption{Características técnicas relevantes ao protótipo}
\label{tab:componentes}
\small
\begin{tabularx}{\textwidth}{L{2.7cm}Y Y}
\toprule
Componente & Características de referência & Uso e limitação no projeto\\
\midrule
Heltec WiFi LoRa 32 V2 &
ESP32 clássico, dois núcleos LX6 até 240~MHz, 520~KiB de SRAM no chip e
8~MB de flash na placa; Wi-Fi de 2,4~GHz. &
Aquisição e inferência local. A memória disponível à aplicação depende
da partição e dos recursos do firmware.\\
DHT22 / AM2302 &
Temperatura de $-40$ a $80$\celsius; UR de 0 a 100\%; alimentação de
referência de 3 a 5~V; novas leituras em intervalos de pelo menos 2~s. &
Alimentado em 3,3~V. A referência consultada informa erro de
$\pm0{,}5$\celsius\ e de 2 a 5 pontos percentuais de UR; não aferidos no protótipo.\\
Termopar tipo K &
Sensor termoelétrico passivo; a conversão depende do MAX6675 e da
compensação da junção fria. &
Mede o ponto em que a ponta está instalada. Não caracteriza, sozinho,
o campo térmico de toda a massa de grãos.\\
MAX6675 &
Conversor de 12 bits, resolução de 0,25\celsius, alimentação de 3,0 a 5,5~V
e conversão de até 220~ms. &
Alimentado em 3,3~V; leitura serial por software. O código aceita
0 a 1024\celsius; o maior código de 12 bits equivale a 1023,75\celsius.\\
Módulo SW-420 &
Saída digital de um módulo de detecção de vibração; o ajuste e a
polaridade dependem da implementação utilizada. &
Alimentado em 3,3~V. Registra DO e transições; não fornece aceleração,
intensidade ou frequência mecânica calibradas.\\
\bottomrule
\end{tabularx}
\source{especificações de referência de Heltec, Espressif, Adafruit e Analog Devices; limites de uso extraídos do firmware \cite{heltecV2,espressif2026,adafruitDHT,max6675}}
\end{table}

No MAX6675, a resolução de 0,25\celsius\ não significa que o conjunto
termopar--conversor tenha esse erro máximo. O documento do fabricante
estabelece, por exemplo, erro do conversor de até 8~LSB na faixa de 0 a
700\celsius, com alimentação de 3,3~V e temperatura ambiente especificada,
além de tratar separadamente a compensação da junção fria
\cite{max6675}. Essas condições não constituem calibração da sonda montada.
Para o SW-420 não foi identificado um certificado ou uma especificação
metrológica do módulo efetivamente utilizado; por isso, não se atribuem
unidades físicas à contagem digital.

\subsection{Periodicidade e janela de vibração}

O firmware estabelece um intervalo nominal de 2500~ms. Na ausência de
atrasos adicionais, a taxa de ciclos é
\begin{equation}
f_{\mathrm{nominal}}=\frac{1}{2{,}5~\mathrm{s}}
=0{,}4~\mathrm{Hz},
\end{equation}
equivalente a aproximadamente 24 ciclos por minuto. Essa é a taxa nominal
de aquisição e envio, não uma frequência de amostragem de uma forma de
onda de vibração. A duração real é transmitida em \texttt{window\_ms},
pois atrasos de rede ou de execução podem prolongar a janela.

As transições do SW-420 são capturadas por uma interrupção em
\texttt{CHANGE}; o contador é copiado e reiniciado em seção crítica.
O fechamento da janela antecede a leitura dos sensores térmicos, portanto
o registro não deve ser descrito como uma aquisição simultânea de todas
as grandezas. A rotina de interrupção não faz HTTP, leitura do DHT22 ou
impressão serial. Como a biblioteca do DHT pode bloquear interrupções
brevemente, a contagem continua sendo uma medida demonstrativa.

A janela inicial também fornece tempo superior aos 220~ms de conversão
do MAX6675. A leitura ocorre uma vez por ciclo, utilizando o construtor
na ordem SCK, CS e SO. Não se pressupõe o uso de I$^{2}$C: o DHT22 usa
seu protocolo temporizado, e o MAX6675 utiliza uma interface serial
compatível com SPI implementada por software.

\subsection{Validação, atributos e tratamento de falhas}

O DHT22 é considerado válido quando temperatura e UR são finitas e estão
nos limites aceitos pelo firmware. Se falhar, ambos os valores são enviados
como \texttt{null}, com \texttt{dht\_ok=false}. A sonda é validada de forma
independente; termopar aberto pode resultar em \texttt{NAN} na biblioteca,
tratado como temperatura \texttt{null} e \texttt{probe\_ok=false}.
Temperatura igual a zero é uma entrada possível e não identifica falha.

Somente com os dois sensores térmicos válidos é preparado o vetor
\begin{equation}
\mathbf{x}=
\left[T_{\mathrm{ar}},\ \mathrm{UR}_{\mathrm{ar}},\
T_{\mathrm{sonda}},\ \Delta T\right],
\qquad
\Delta T=T_{\mathrm{sonda}}-T_{\mathrm{ar}}.
\label{eq:atributos}
\end{equation}
O pré-processamento reproduz operações em \texttt{float32}. A versão
embarcada não calcula médias móveis, espectro, tendências ou dias de
armazenamento como entradas da IA. Também não implementa filtragem
digital de vibração. Essas operações não devem ser atribuídas ao
protótipo apenas por aparecerem em propostas de evolução.

Uma falha térmica torna a classificação indisponível, sem reaproveitar
uma classe anterior como se fosse atual. A telemetria de vibração
continua sendo enviada. De modo recíproco, uma leitura de vibração
igual a zero não impede a inferência térmica.
\FloatBarrier
